% Vietnamese translation for OI Public Library
% Source-PDF: IOI中国国家候选队论文/1999/杨帆.pdf
\documentclass[11pt]{article}
\usepackage{oipl-vn}
\usepackage{graphicx}
\renewcommand{\figurename}{Hình}

\title{Đúng, đủ, đẹp: ba yếu tố trong thiết kế dữ liệu kiểm thử}
\author{Dương Phàm}
\date{}

\begin{document}
\maketitle

\origtitle{Accuracy, completeness, and aesthetics: three elements of test-data design}
\sourcepath{IOI China candidate team papers / 1999 / Yang Fan.pdf}

\paragraph{Từ khóa.}
Dữ liệu kiểm thử; tính chính xác; tính toàn diện; tính thẩm mỹ.

\section*{Tóm tắt}

Dữ liệu kiểm thử là một bộ phận hữu cơ không thể thiếu của các kỳ thi tin học
hiện đại, và vai trò của nó ngày càng được coi trọng. Bài viết này đề xuất
cách nhìn ``tín, đạt, nhã'' trong thiết kế dữ liệu kiểm thử: chính xác, toàn
diện và đẹp; sau đó lần lượt bàn về ba yếu tố ấy.

\section{Dẫn nhập}

Từ khi Olympic Tin học Quốc tế (IOI) bắt đầu năm 1989, kỳ thi này đã đi qua
mười mùa; Olympic Tin học Toàn quốc Trung Quốc (NOI) cũng đã có lịch sử mười
lăm năm. Trong quá trình phát triển đó, các kỳ thi tin học liên tục tự hoàn
thiện và hình thành một phương pháp riêng. Một điểm khác biệt lớn so với nhiều
môn khác là cách chấm điểm. Các môn truyền thống thường cho điểm từng bước,
tức mỗi bước trung gian có số điểm tương ứng. Còn trong tin học, ngay từ các
kỳ IOI đầu tiên đã dùng kiểm thử hộp đen: không đọc và phân tích mã nguồn, mà
chỉ dựa trên tính đúng đắn của kết quả chương trình với dữ liệu kiểm thử đã
cho; toàn bộ quá trình do máy tính thực hiện tự động.

So với kiểm thử hộp trắng trước đây, kiểm thử hộp đen khách quan, công bằng và
hiệu quả hơn, nên hiện nay được dùng rộng rãi. Trong quá trình ấy, bộ dữ liệu
kiểm thử giữ vai trò cực kỳ quan trọng, thậm chí không kém bản thân đề bài.
Một đề rất hay nhưng không có dữ liệu tốt thì giá trị bị giảm mạnh. Ngược lại,
một đề tưởng như bình thường, nếu đi kèm một bộ dữ liệu xuất sắc, có thể tạo
ra hiệu quả đáng kinh ngạc. Dữ liệu kiểm thử cũng là một cửa ải quan trọng để
điều chỉnh độ khó: cùng một đề, dùng các bộ dữ liệu khác nhau có thể dẫn tới
độ khó rất khác nhau. Vì vậy dữ liệu kiểm thử là một phần không thể thiếu của
bài thi tin học, và việc thảo luận về nó là cần thiết.

Học giả Nghiêm Phục từng nêu tiêu chuẩn ``tín, đạt, nhã'' cho dịch thuật:
tín là chân thực và xác đáng; đạt là thông suốt, có thể hiểu rộng là đầy đủ;
nhã là cao đẹp. Dù dịch thuật văn học và thi tin học không có liên hệ trực
tiếp, ba chữ đó vẫn có thể vận dụng cho thiết kế dữ liệu kiểm thử. Tín tương
ứng với tính chính xác, đạt tương ứng với tính toàn diện, và nhã tương ứng với
tính thẩm mỹ. Một bộ dữ liệu xuất sắc phải là sự kết hợp tốt của cả ba.

\section{Tính chính xác}

Tầm quan trọng của tính chính xác trong dữ liệu kiểm thử là tuyệt đối. Mất đi
tính chính xác thì mọi khía cạnh khác đều không còn cơ sở. Với một bài thi tin
học, nếu dữ liệu kiểm thử chứa lỗi thì giá trị thi đấu của bài đó bằng không.
Cụ thể, tính chính xác thể hiện ở ba mặt: định dạng vào ra phù hợp yêu cầu đề;
dữ liệu nằm trong phạm vi ràng buộc của đề; và đáp án chuẩn đúng. Thiếu bất kỳ
mặt nào cũng không được.

Trước hết, định dạng vào ra phải khớp với đề bài. Sau khi chuyển sang chấm tự
động, tầm quan trọng của định dạng trở nên rất rõ: lỗi định dạng và lỗi thuật
toán đều có cùng hậu quả. Máy chấm chỉ so sánh tệp xuất của thí sinh với tệp
chuẩn; khác nhau là bị chấm sai. Vì vậy nếu tệp xuất chuẩn trong dữ liệu bị
sai, mọi lời giải đúng của thí sinh cũng có thể bị đánh sai. Tương tự, nếu
định dạng đầu vào sai, hoặc không sắp xếp như đề yêu cầu, chương trình đúng có
thể không đọc được dữ liệu đúng và sinh kết quả sai. Tai nạn như vậy trong kỳ
thi là rất nghiêm trọng.

Thứ hai, dữ liệu nhất định phải nằm trong phạm vi giới hạn của đề. Phạm vi này
có nhiều nghĩa: quy mô dữ liệu, cận trên và cận dưới, kiểu dữ liệu, v.v. Ví
dụ, trong bài IOI 1998 ``Contact'', giới hạn kích thước tệp vào tới $2$ MB là
giới hạn quy mô; các điều kiện $0<n\le 20$ và $0<a\le b\le 12$ là cận của dữ
liệu; còn ``dãy nhị phân kết thúc bởi ký hiệu 2'' là kiểu dữ liệu. Dữ liệu
kiểm thử không được vượt qua bất kỳ ranh giới nào.

Trong luyện tập, người ta thường tạo dữ liệu rất lớn để đo tốc độ chương
trình, nhất là với bài tìm kiếm. Cách kiểm thử này ẩn chứa nguy cơ lớn: quy mô
có thể vượt giới hạn đề, hoặc nhìn bề ngoài vẫn hợp lệ nhưng trong quá trình
tính toán làm phát sinh tràn giá trị trung gian. Đây là loại lỗi dễ xảy ra và
khó phát hiện. Chẳng hạn bài IOI 1998 ``Polygon'' quy định tọa độ đỉnh luôn
nằm trong đoạn $[-32768,32767]$; nhiều dữ liệu có quy mô lớn có thể cho kết
quả cuối cùng trong đoạn đó nhưng giá trị trung gian vượt ngoài đoạn, và như
vậy không còn là dữ liệu hợp lệ.

Cuối cùng, và quan trọng nhất, kết quả chuẩn phải đúng. Điều này không chỉ là
yêu cầu của bài thi tin học mà còn là yêu cầu chung của mọi bài thi. Để bảo
đảm điều đó trong tin học, trước hết phải bảo đảm chương trình chuẩn đúng, vì
kết quả chuẩn được sinh từ chương trình chuẩn. Lỗi trong chương trình chuẩn sẽ
dẫn trực tiếp tới lỗi trong dữ liệu kiểm thử. Khi thiết kế dữ liệu phải đặc
biệt chú ý điểm này.

Tóm lại, chỉ khi thỏa ba mặt trên, dữ liệu kiểm thử mới bảo đảm được tính đúng
cần có. Dữ liệu đúng mới là dữ liệu đạt yêu cầu tối thiểu.

\section{Tính toàn diện}

Nghiêm Phục nói: nếu chỉ tín mà không đạt thì tuy gọi là dịch cũng như chưa
dịch. Với dữ liệu kiểm thử cũng vậy. Chỉ có chính xác thì chưa thể gọi là dữ
liệu tốt; đó chỉ là yêu cầu tối thiểu. Sau khi bảo đảm chính xác, phải xét tới
tính toàn diện, bởi chỉ như vậy mới phân biệt được chất lượng các chương trình
và đạt mục đích của kỳ thi.

Tính toàn diện của dữ liệu kiểm thử chủ yếu thể hiện ở hai mặt: kiểm tra các
trường hợp đặc biệt, và kiểm tra hiệu suất thuật toán cũng như khả năng chịu
đựng về thời gian, bộ nhớ của chương trình.

\subsection{Trường hợp đặc biệt}

Trường hợp đặc biệt lại có hai loại. Một là điều kiện biên của đề; hai là
những tình huống không bị đề cấm rõ ràng nhưng trái với trực giác thông
thường. Cả hai đều dễ bị bỏ sót, đặc biệt là loại thứ hai.

Mỗi bài đều có điều kiện biên. Khi dữ liệu vượt qua biên, nó mất tính chính
xác. Nhưng khi dữ liệu nằm đúng trên biên của đề, nó đạt mục đích kiểm tra
biên. Dữ liệu như vậy vừa chính xác vừa có giá trị. Thông thường bài toán có
cận dưới và cận trên. Kiểm tra cận dưới thường thuộc nhóm trường hợp đặc biệt,
còn kiểm tra cận trên thường cần quy mô lớn và liên quan tới hiệu suất. Ví dụ
bài IOI 1998 ``Starry Night'' có giới hạn $0\le$ số chòm sao $\le 500$; ta có
thể thiết kế bản đồ sao rỗng để kiểm tra trường hợp số chòm sao bằng $0$.
Tương tự, trong bài ``Camelot'', nếu số hiệp sĩ có thể bằng $0$, dữ liệu không
có hiệp sĩ sẽ kiểm tra mức độ toàn diện trong suy nghĩ của thí sinh.

Một câu nói quen thuộc trong khoa học là: điều gì không bị cấm tuyệt đối thì
đều có thể xảy ra. Trong thi tin học cũng vậy. Miễn là đề không nói rõ cấm,
tình huống đó đều có khả năng và có lý do xuất hiện trong dữ liệu. Chẳng hạn
một dữ liệu của bài NOI 1997 ``File Matching'' có trường hợp cùng một tệp vừa
bị yêu cầu thao tác vừa bị yêu cầu không thao tác. Vì đề không cấm rõ điều
này, và dữ liệu kiểm tra đúng điểm đó, nó là một dữ liệu đẹp. Các chi tiết như
vậy vừa dễ bị thí sinh bỏ qua, vừa dễ bị người ra đề bỏ qua, nên cần được coi
trọng.

\subsection{Hiệu suất và sức chịu đựng}

Kiểm tra hiệu suất thuật toán và khả năng chịu đựng thời gian, bộ nhớ cũng rất
quan trọng. Với cùng một bài, các thí sinh có thể xây dựng mô hình toán học và
thuật toán khác nhau, trong đó nhiều thuật toán đều có vẻ khả thi. Khi chấm,
cần dùng dữ liệu để phân biệt hiệu quả của các thuật toán ấy.

Cách thông thường để kiểm tra hiệu suất là tạo dữ liệu quy mô lớn. Ví dụ với
bài ``Polygon'', khi quy mô tăng, lời giải tìm kiếm có thời gian chạy tăng rất
nhanh, còn lời giải quy hoạch động ổn định hơn. Dữ liệu như vậy phân biệt rõ
thuật toán tốt và kém.

Ngay cả khi hai thí sinh dùng cùng một thuật toán, cách cài đặt và cấu trúc dữ
liệu cũng có thể khác nhau, nên cần kiểm tra khả năng chịu đựng của chương
trình. Khả năng chịu về thời gian không chỉ phụ thuộc độ phức tạp thuật toán,
mà còn phụ thuộc tối ưu hóa, tiền xử lý và nhiều chi tiết khác; điều này đặc
biệt rõ trong các bài tìm kiếm. Để phân biệt các chương trình cùng bậc độ
phức tạp, có thể thiết kế dữ liệu mà chương trình có tối ưu hoặc tiền xử lý
chạy được trong thời hạn, còn chương trình không có thì không. Ví dụ với bài
``Picture'', dữ liệu hình chữ nhật chứa hoàn toàn một hình khác sẽ tạo lợi thế
rõ cho chương trình có tiền xử lý phù hợp.

Khả năng chịu về bộ nhớ lại gắn chặt với cấu trúc dữ liệu. Một cấu trúc tốt
có thể tiết kiệm rất nhiều bộ nhớ, nhưng điều này không phải lúc nào cũng thể
hiện qua thời gian chạy. Vì vậy cần tạo dữ liệu để phân biệt lựa chọn cấu trúc
dữ liệu: với bài tìm kiếm có thể tăng độ sâu tìm kiếm; với bài xử lý lượng dữ
liệu lớn có thể tăng dung lượng; với bài tính toán nặng có thể tăng quy mô.
Khi đó chương trình dùng cấu trúc không phù hợp có thể tràn ngăn xếp, tràn
vùng nhớ, vượt phạm vi lưu trữ, hoặc thậm chí treo. Chỉ những chương trình vừa
có thuật toán hiệu quả vừa chọn cấu trúc dữ liệu thích hợp mới vượt qua.

Do đó, để kiểm thử toàn diện, không thể máy móc quy định mọi bài đều có $5$
hay $10$ bộ dữ liệu. Số lượng tốt nhất phụ thuộc vào bản thân bài toán, độ khó
và quy mô yêu cầu. Ít quá thì không toàn diện; nhiều quá thì tác dụng bổ sung
nhỏ. Một bài đơn giản và quy mô nhỏ như ``Camelot'' có thể chỉ cần ít dữ liệu;
trong khi bài ``Picture'' có quy mô lớn, nhiều tình huống và nhiều thuật toán
khác nhau thì cần nhiều dữ liệu hơn.

Độ khó của dữ liệu kiểm thử thường quyết định độ khó của bài. Với một kỳ thi
thông thường, có thể phân bố điểm kiểm thử như sau:
\[
\begin{array}{l|r}
\text{Dữ liệu đơn giản, quy mô nhỏ} & 25\%\\
\text{Trường hợp đặc biệt} & 20\%\\
\text{Kiểm tra hiệu suất thuật toán, quy mô lớn} & 20\%\\
\text{Kiểm tra sức chịu thời gian và bộ nhớ, quy mô lớn} & 35\%
\end{array}
\]

Bốn nhóm này không tách rời tuyệt đối, vì dữ liệu kiểm tra bộ nhớ và thời gian
cũng thường kiểm tra hiệu suất thuật toán. Với kỳ thi ở mức thấp hơn hoặc cao
hơn, tỉ lệ nên thay đổi:
\[
\begin{array}{l|rr}
 & \text{Dễ hơn} & \text{Khó hơn}\\
\hline
\text{Dữ liệu đơn giản, quy mô nhỏ} & 55\% & 5\%\\
\text{Trường hợp đặc biệt} & 10\% & 15\%\\
\text{Kiểm tra hiệu suất thuật toán, quy mô lớn} & 20\% & 25\%\\
\text{Kiểm tra sức chịu thời gian và bộ nhớ, quy mô lớn} & 15\% & 55\%
\end{array}
\]

Với cùng một đề, tỉ lệ thứ nhất làm bài dễ hơn, tỉ lệ thứ hai làm bài khó hơn.
Như vậy dữ liệu kiểm thử có chức năng điều chỉnh độ khó. Một ví dụ điển hình
là bài đầu tiên của NOI 1998, ``Personal Income Tax''. Bản thân bài không khó,
thuật toán gần như hiển nhiên. Nhưng nó xuất hiện ở vòng toàn quốc và tỉ lệ
đạt điểm chỉ khoảng $11\%$, vì dữ liệu kiểm thử đã đóng vai trò điều chỉnh:
trong $10$ bộ dữ liệu chuẩn, có tới $9$ bộ kiểm tra khả năng chịu đựng của
chương trình, và chỉ dùng tính toán độ chính xác cao mới ra đáp án đúng. Nếu
khi đó dùng dữ liệu quy mô nhỏ, tỉ lệ đạt điểm có lẽ vượt $90\%$ và bài chỉ
còn là bài dễ ở vòng khu vực.

Vì thế tỉ lệ độ khó của dữ liệu kiểm thử cực kỳ quan trọng. Nhưng tất cả điều
đó đều dựa trên dữ liệu toàn diện; nếu thiếu tính toàn diện, dữ liệu khó tác
động chính xác tới độ khó của bài. Tính toàn diện vì vậy là điểm then chốt.

\section{Tính thẩm mỹ}

Chỉ chính xác và toàn diện đã đủ để gọi là dữ liệu tốt chưa? Câu trả lời là
chưa. Với dữ liệu kiểm thử cũng cần ``ngoài tín và đạt, còn cầu nhã''. Ở đây
``nhã'' không phải sự trau chuốt văn chương, vì dữ liệu chỉ gồm ký tự và số,
khó có thể sinh động theo nghĩa văn học. Cái đẹp của dữ liệu kiểm thử có ý
nghĩa riêng, nhưng lâu nay người ra đề chưa thật sự coi trọng.

Có câu: chuẩn mực là đẹp, hài hòa là đẹp. Cái đẹp của dữ liệu kiểm thử được
xây trên hai điểm đó. Dữ liệu không chỉ phục vụ máy tính chấm tự động; cuối
cùng nó vẫn được con người đọc. Sau kỳ thi, thí sinh hoặc huấn luyện viên khi
phân tích, tổng kết thường phải xem dữ liệu. Lúc ấy, dữ liệu hài hòa và có
quy phạm giúp ích hơn rất nhiều so với dữ liệu lộn xộn.

Vậy thế nào là dữ liệu chuẩn mực và hài hòa? Trong bài gốc, tác giả lấy bài
``Picture'' của IOI 1998 và đưa ra năm bộ dữ liệu được chuyển thành hình trực
quan. Một số hình tạo cảm giác rõ ràng, cân đối hơn hẳn các hình khác; hình
lộn xộn khiến người đọc khó nhận ra cấu trúc, dù đã được vẽ lại trực quan,
huống hồ khi chỉ nhìn dữ liệu thô. Qua đó có thể thấy chuẩn mực và hài hòa là
khái niệm tương đối, nhưng nếu chỉ nhìn dữ liệu mà người đọc có thể hình dung
ngay một hình ảnh trực quan trong đầu, dữ liệu ấy có thể gọi là đẹp.

\begin{figure}[ht]
\centering
\includegraphics[width=0.92\linewidth]{assets/ioi/1999/yang-fan-picture-figures.png}
\caption{Năm bộ dữ liệu của bài IOI 1998 ``Picture'' sau khi chuyển thành hình trực quan trong bản gốc.}
\end{figure}

Trong năm hình này, Hình 1 và Hình 4 cho cảm giác tốt hơn nhiều so với Hình 2
và Hình 3, còn Hình 5 nằm giữa hai nhóm. Đó chính là tác dụng của sự chuẩn mực
và hài hòa: Hình 2 và Hình 3 đã gây cảm giác rối ngay cả khi được biểu diễn
thành hình ảnh, nên khi chỉ đọc dữ liệu thô còn khó nhận ra cấu trúc hơn.

Theo tiêu chuẩn này, trong dữ liệu của nhiều kỳ thi trước đây, những bộ thật
sự ``nhã'' là rất hiếm; phần nhiều là dữ liệu sinh ngẫu nhiên. Không thể phủ
nhận dữ liệu ngẫu nhiên có tính khái quát, và trong mỗi bài thường cần trộn
một số dữ liệu ngẫu nhiên để đạt tính toàn diện. Nhưng nếu lạm dụng, dữ liệu
ngẫu nhiên không chỉ không cần thiết mà còn phá vỡ cảm giác đẹp của cả bài.
Hơn nữa, khi thí sinh phân tích sau kỳ thi, nếu gặp toàn dữ liệu ngẫu nhiên,
họ khó hiểu dữ liệu và khó phân tích chương trình. Nếu dữ liệu đẹp, người đọc
có thể phối hợp trực giác của mình với kết quả máy tính để phân tích sâu hơn,
hiệu quả hơn.

Vì vậy tính thẩm mỹ của dữ liệu kiểm thử cũng rất cần thiết. Giá trị của một
bài thi không chỉ nằm trong lúc thi, mà còn nằm ở việc sau đó thí sinh phân
tích, tổng kết và nâng cao trình độ. Dữ liệu đẹp giúp phát huy giá trị sau kỳ
thi của bài toán. Trong các kỳ thi tương lai, mức độ coi trọng điểm này chắc
chắn sẽ tăng lên.

\section{Kết luận}

Tóm lại, tín, đạt, nhã là ba yếu tố của thiết kế dữ liệu kiểm thử. Chỉ khi có
tính chính xác, tính toàn diện và tính thẩm mỹ, dữ liệu mới có thể gọi là xuất
sắc. Tuy nhiên tầm quan trọng của ba yếu tố không hoàn toàn giống nhau: tín là
nền tảng, đạt là then chốt, nhã là nâng cao. Không có tín thì dữ liệu mất cơ
sở tồn tại, đạt và nhã không còn ý nghĩa. Chỉ chính xác chưa đủ; phải toàn
diện, và tính toàn diện thường cần cả một nhóm dữ liệu mới đạt được. Một bài
hay không chỉ có giá trị trong cuộc thi, mà còn giúp thí sinh học được sau khi
phân tích; trong quá trình ấy, dữ liệu đẹp phát huy tối đa tiềm năng của đề.

Một bộ dữ liệu xuất sắc có ảnh hưởng không thể xem nhẹ: nó có thể trực tiếp
quyết định độ khó, và đôi khi khiến một bài tưởng như bình thường tạo ra hiệu
quả bất ngờ. Vì thế dữ liệu kiểm thử là bộ phận không thể thiếu của bài thi
tin học, cần được coi trọng và nghiên cứu sâu hơn. Những ý kiến trên chỉ là
cách nhìn cá nhân của tác giả, mong gợi mở thêm thảo luận giữa thí sinh và
huấn luyện viên về chủ đề thiết kế dữ liệu kiểm thử.

\end{document}
